\subsection{The Case II starting tower and marked rational lattice}
\label{subsec:rh-four-rational-lattice}

We give the reduction for the actual remaining RH datum of
Theorem~\ref{thm:rh-symplectic-arithmetic-data}. In particular
$\ell\nmid N$, $\psi\ne1$, $\psi$ has power-of-two order,
$\psi$ is unramified at $\ell$, and $\mathrm{Sb}=\varnothing$.
These simplifications are specific to the RH host. Write
$x=\bar\xi$, $y=\bar\Psi=\bar\epsilon^e\bar\psi$ and
$\omega=\bar\epsilon$. We use the original two root coordinates:
\[
 A=w_{\rm Al}=(1,0),\quad P=w_{\rm Ps}=(0,1),
 \quad B=w_{\rm Be}=(-1,1),\quad H=w_{\rm Th}=(1,1).
\]
The coordinate weights of the four basis vectors are respectively
$(1,1),(0,1),(1,0),(0,0)$, and the weight of entry $ij$ is their
difference. The symplectic form pairs positions $0,3$ and $1,2$,
with the signs in the original adjoint.

\begin{lemma}[The original tame fitted tower for RH Case II]
\label{lem:rh-four-tame-start}
For either permissible labeling $\mathrm{pr}$ of the original
datum, there is a tower $T_0$ satisfying $\mathrm{Fit}$,
$\mathrm{Two}$, $\mathrm{TameL}$ and $\mathrm{Deep}$.
Its initial fields can be chosen unramified outside $S$ and three
new rational primes avoiding any prescribed finite set.
\end{lemma}

\begin{proof}
Choose three quadratic characters $\phi,\eta,\eta'$ of distinct
prime conductors outside $S$ and the prescribed finite set.
Take $\phi$ even and $\eta,\eta'$ odd, and require each to be
trivial on the full decomposition groups at the primes of $S$.
This is achieved by primes respectively congruent to $1,-1,-1$
modulo $8\prod_{q\in S}q$, by Dirichlet's theorem and quadratic
reciprocity (for the two odd characters use the discriminant
$-p$). Define
\[
 G_F=\{g:y(g)=\phi(g),\quad\eta(g)=\eta'(g)\},\qquad
 K=\overline{\mathbb Q}^{\ker\eta},\qquad L=FK.
\]
The kernel conditions give a finite abelian extension $F$.
It is real, since $y$ is even, $\phi$ even and both other
characters odd. Its degree is even because $\eta\eta'$ is a
nontrivial quadratic character trivial on $G_F$. The field
$L$ is CM with maximal real field $F$. On inertia in $G_F$,
the common value of $\eta,\eta'$ is one, since their conductors
are disjoint. Thus $L/F$ is everywhere unramified at finite
places. All primes of $S$ split in $K$. The group index of
$G_F$ is prime to $\ell$, so the original $\mathrm{TameL}$
condition holds, since the residual character $y$ and the
three quadratic characters have order prime to $\ell$.
Since $\psi$ has order prime to $\ell$, $\mathrm{kc}$ and
$\mathrm{Deep}$ hold automatically. For $\mathrm{kb}$, away
from $\ell$ the cyclotomic character is unramified, so that
triviality of the full residual diagonal on inertia implies
$\bar\psi=1$, hence $\psi=1$. On $G_F$ the character $y$ is
quadratic. It is unramified: away from the conductor of $\phi$
use $y=\phi$, and at that conductor use that $y$ itself is
unramified. It is trivial on the relative $\ell$-adic
decomposition groups, as $\phi$ is split there. Since $\xi$
is unramified at $\ell$, the whole residual diagonal is
trivial on relative inertia. The real quadratic inducing
field of $\mathrm{Two.quad}$ is the field of $\phi$, split
at $2,\ell$. Its restriction is nontrivial by the witnesses
below. The labeling sign replaces $y$ by its inverse in
root formulas and therefore has no effect on these assertions.

Here are explicit witnesses for all remaining conditions,
including the order-four conditions. Write $b$ for the dyadic
conductor exponent of $\psi$, and $2^n$ for the original
conductor of $\xi$. The order of $\psi$ divides $2^{b+2}$.
At $\ell$, take the original Frobenius lift $s_\ell$ with
$\epsilon(s_\ell)=1$, and raise it to $2^{b+2}$.
Then $y=1$ and $\omega=1$. The value of $x^4$ is not one:
otherwise faithfulness of $\xi$ on $1+4\mathbb Z/2^n$
would give
$\ell^{2^{b+4}}\equiv1\bmod 2^n$, contradicting
\[
 n\ge b+6+3\log_2\ell>
 v_2(\ell^2-1)+b+3 .
\]
Here we used the elementary valuation identity
$v_2(\ell^{2^j}-1)=v_2(\ell^2-1)+j-1$.
The witness already belongs to $G_F$ and to $G_K$,
as the three new characters split at $\ell$.
At $2$, take an inertia element with cyclotomic image
$3\bmod2^n$ in the dyadic cyclotomic field, and raise it
to $2^b$. The dyadic character $\psi_2$ has order dividing
$2^b$, the odd character and $\epsilon$ are unramified
there, and $x(3)$ has order $2^{n-2}$: the square of $3$
generates the subgroup of order $2^{n-3}$ of the principal
units, on which $\xi$ is faithful. Since $n>b+4$, the
result has $x^4\ne1$, $y=\omega=1$. Again it lies in
$G_L$. This proves $\mathrm{bigL},\mathrm{big2}$ and
$\mathrm{locL}$, the latter by its second alternative.

Multiplication by suitable inertia elements at the three
new primes can independently change the three quadratic
signs without changing $x,y,\omega$; these original characters
are unramified there. Thus we can move any prescribed element
with $y=\pm1$ into $G_L$ with the same values of $x,y,\omega$.
Moving complex conjugation gives $a\in G_L$ with
$(x,y,\omega)=(-1,1,-1)$. Let $u$ be the preceding
witness at $2$, so $x(u)^4\ne1$, $y(u)=\omega(u)=1$.
Since $\psi$ is a nontrivial power-of-two order character
of conductor prime to $\ell$, there is an element with
$\psi=-1$ and $\epsilon=1$: its finite cyclotomic field
is disjoint from the $\ell$-power cyclotomic field.
Move it into $G_L$ by the sign changes, and if necessary
multiply by $u$. This gives $t\in G_L$ with
$y(t)=-1,\ \omega(t)=1,\ x(t)^2\ne1$.
Take the mark $g_0=t$. The four diagonal values
$x(t)y(t),y(t),x(t),1$ are distinct. Every nonzero root
has nontrivial value at $u$ or $t$, and $u$ gives the
required order greater than four for $A$.
For every root-product character $x^i y^{\pm j}$,
if $i$ is even it differs from $\omega$ at $a$, and if
$i$ is odd it differs at $u$, because $x(u)$ has power-of-two
order greater than four. This proves $\mathrm{g2}$ for
all integer pairs. A quadratic character of $G_L$ permuting
the diagonal characters must be one of $1,x,y,xy$, by its
action on $1$. If trivial on $\ker\omega$ it must be $1$,
using $u,t$. This verifies $\mathrm{g3'}$ exactly as stated,
including its character-valued condition rather than merely
an exclusion of one particular quadratic field.

Choose for each prime of $F$ over $\ell$ an absolute
place extending the fixed place of $K$ above $\ell$.
This is possible since $K$ is split at $\ell$ and
$F\cap K=\mathbb Q$; the action of the absolute Galois
group over $F$ can interchange the two places of $K$.
The usual transitive action on places extending a fixed
finite place supplies the original \texttt{cover} and
\texttt{one}. All fields of $\mathrm{Ok}$ and
$\mathrm{Two}$ have now been verified with the specified
labeling and mark.
\end{proof}

\begin{theorem}[Rational marked lattice reduction]
\label{thm:rh-four-lat-shape}
Let $(r,u)$ be an original \texttt{KPair} in any of its
ambient valued algebraically closed fields, and suppose its
fraction-field representation satisfies \texttt{Star}.
There are a permissible labeling $\mathrm{pr}$ and a
tower as in Lemma~\ref{lem:rh-four-tame-start} such that the
original rational \texttt{QPair} admits a conjugate integral
model with \texttt{LatShape} at its mark. The ordered
inertial characters remain the original $\mathrm{ps}_i$,
not a relabeled list. The resulting pair still satisfies
the original \texttt{Small} condition.
\end{theorem}

\begin{proof}
Work over the complete DVR inside the ambient valuation
ring supplied by $\mathrm{Small}$, enlarging its fraction
field by finite algebraic extensions when needed. Its
residue field is finite. The residue map into the ambient
residue field is injective and the ambient field is algebraic
over its fraction field. A finite extension of a complete
discretely valued field has a unique extended valuation,
complete DVR and finite residue field; the given embedding
therefore preserves integrality and continuity. All matrix
operations below take place in such rings.

First choose the labeling. On the residual space the
original $\ell$-adic marked element $s_\ell$ separates
the blocks $\{0,2\},\{1,3\}$ by \texttt{condD}.
The first column of the ordinary flag matrix is a
nonzero eigenvector in the first block. Its component
at $0$ or $2$ is nonzero. In the latter case if needed
apply the symplectic permutation interchanging $0,2$
and $1,3$ (with the signs preserving the form), so the
first component is a unit. The interchange is permitted:
if the original component at $2$ is nonzero, the
first-column identity for the flag for every local group
element implies $\bar{\mathrm{ps}}_0=\bar{\mathrm{ps}}_2$,
hence $y=1$ on the entire decomposition group. This is
precisely \texttt{LabOk}. The first two columns modulo
the maximal ideal span a stable Lagrangian. Their span
splits by the two separated blocks at $s_\ell$; the
second graded character belongs to the other block.
Thus the intersection with the first block is just
the first line, and with the second block a line.
By symplectic orthogonality that second line has
nonzero component at $1$ when the first line has
nonzero component at $0$: explicitly if the first
line has coordinates $(a,0,c,0)$ and the second
$(0,b,0,d)$, the equation is $ad=bc$ with the
source signs, and $a\ne0$ precludes $b=0$.
Consequently the first two leading minors of the
relabeled flag matrix are units.

Choose the tower of Lemma~\ref{lem:rh-four-tame-start}
for this labeling. Its mark has pairwise distinct residual
eigenvalues. Hensel's lemma gives the four uniquely labeled
eigenvalues over the complete DVR, with unit differences,
and their spectral projectors give eigenvectors congruent
to the coordinate basis. The fixed similitude pairs the
complementary eigenvalues: the form vanishes between
eigenspaces unless the product of their eigenvalues is
the multiplier, and the complementary pairings are units
by reduction. Rescale two eigenvectors by these units to
obtain a symplectic diagonalization congruent to the
identity. If the original representation is block
parabolic in $\{0,2\}\mid\{1,3\}$, the relabeling
preserves this parabolic, as does the diagonalization:
the eigenprojectors are polynomials in the marked matrix,
so eigenvectors obtained from columns respect the invariant
block.

The transformed flag matrix still has unit leading minors
of sizes one and two. For a symplectic matrix these imply
the unit third leading minor, by complementary minors
and inverse-transpose symmetry. Gaussian elimination
therefore gives a unique lower-unitriangular times upper
triangular factorization. Both factors are symplectic:
apply the group automorphism $g\mapsto J^{-1}g^{-t}J$
for the antidiagonal alternating form $J$, which preserves
each triangular subgroup, and use uniqueness.
Take the lower factor $n$. This leaves every ordered
local diagonal character unchanged. By the preceding
residual block description its entries $n_{10},n_{30},
n_{21}$ vanish modulo the maximal ideal. Symplecticity
gives $n_{32}=-n_{10}$ and
$n_{31}=n_{20}-n_{21}n_{10}$ in the source signs,
so only the weight $-P$ can remain in the lower flag.
We now have exactly \texttt{NearSplit} and \texttt{QPair}.
The representation entries off the diagonal have positive
valuation.

The original $\mathrm{Star}$ hypothesis gives nonzero
global entries of both weights $\pm P$, and at least one
of $B,-A,-H$ (\texttt{StarE}). Indeed an absolutely
irreducible representation spans the full matrix algebra
over the algebraically closed field. In the block-extension
alternative the quotient block on $\{1,3\}$ spans its full
matrix algebra and therefore gives $\pm P$. Block parabolic
conjugation preserves nonsplitting: if the transformed
representation were block diagonal, removing the block
diagonal part of the conjugating matrix would exhibit the
original splitting by a block unitriangular matrix.
The nonzero cross-block entry has weight in $\{B,-A,-H\}$.

Let $k_w$ be the minimum of the valuations of nonzero
global entries of weight $w$ and of the independent lower
flag entries of that weight, omitting a weight if all are
zero. These minima exist in the value group of the
complete DVR, in units of its uniformizer valuation.
Only $k_{-P}$ may be zero. Consider the rational polyhedron
\[
              w_1 a+w_2 b+k_w\ge0
                       \qquad(k_w\text{ defined}).
\]
It contains the origin. Maximize $b$ and then $a$.
The first maximum is finite by the nonzero $-P$ entry,
and the second is finite by the nonzero entry of weight
$B,-A$ or $-H$. Rational linear programming gives an
attained rational optimum (in fact the denominators divide
two, from the determinants of the root coefficient
vectors). Thus a finite extension containing the appropriate
uniformizer root realizes the optimum by the diagonal
similitude $\mathrm{diag}(t_1t_2,t_2,t_1,1)$ with
valuations $a,b$. Conjugate both the representation and
the flag by this matrix. All entries are integral, using
the two displayed relations for the dependent lower
flag entries; every labeled local character and diagonal
entry is unchanged. Multiplication by the fixed scalars
preserves continuity in the enlarged complete DVR.

Call a root tight if some corresponding representation
entry or independent lower flag entry is nonzero after
reduction. Tight roots are exactly the equalities in the
polyhedron. Since $b\ge0$, $P$ is not tight. Every tight
root other than $-P$ has strictly negative value at
$(a,b)$, since its original valuation was positive.

We can now bypass the six-case enumeration of the original
lattice argument. The odd strict cocycle vanishing already
proved for the original six characters has the following
entry consequence. Suppose all two-step products
$r(g)_{ih}r(g')_{hj}$ vanish for $h\ne i,j$ in the residual
pair. Then $r(g)_{ij}/\bar\chi_j(g)$ is a cocycle of weight
$w_{ij}$ by the multiplication identity and the diagonal
characters. It satisfies the killing and normalization
conditions. If its local entry is
$m(\bar\chi_j-\bar\chi_i)$, the cocycle is a local
coboundary. For odd $w_{ij}$, strict vanishing makes it
a global coboundary; normalization at the distinct mark
makes it zero. The local coefficient $m$ is also zero
because the two characters differ on the $\ell$-adic
decomposition group by \texttt{condD}. This argument
works over any of the residue fields by the finite-span
coefficient argument in the tower-bound theorem.

Two consequences follow from the first-column flag identity.
If $-A$ is not tight, then $r_{10}=r_{32}=n_{10}=0$.
For entry $30$ the two-step products vanish and the
flag identity gives $r_{30}=n_{30}(\bar\chi_0-\bar\chi_3)$
locally. Thus the odd-entry consequence kills $r_{30}$
and $n_{30}$, so $-H$ is not tight. If $B$ and $P$ are
not tight, then $r_{12}=r_{13}=0$. Entry $10$ has no
two-step products, and the first-column identity gives
$r_{10}=n_{10}(\bar\chi_0-\bar\chi_1)$ locally.
It kills $r_{10},n_{10}$, and the symplectic inverse
identity kills $r_{32}$ for all group elements.
Thus $-A$ is not tight.

At the optimum, moving in direction $(1,0)$ must
be blocked by a tight root in $\{-A,B,-H\}$ (otherwise a
small positive step is feasible, with only finitely many
inequalities). If $B$ were not tight, the preceding two
consequences, with $P$ absent, would exclude both remaining
roots. Therefore $B$ is tight. Its equality implies
$a>b\ge0$. Now all of $A,-B,H,P$ have nonnegative scalar
product with $(a,b)$ and positive original valuations,
so none is tight. Only $B,-A,-P,-H$ remain. Among these
roots direction $(-1,1)$ has negative scalar product
only with $-P$. Maximality of $b$ therefore forces
$-P$ tight. This proves exactly the needed root support
$nP$, with $B,-P$ tight, without any sign-pattern case split.

The flag support is consequently $qQ$ and $B$ is a
global nonzero entry, since there is no lower flag
coordinate of weight $B$. If $-P$ is global nonzero
we have the first alternative of \texttt{LatShape}.
Otherwise the global support is $vV$, and $n_{20}\ne0$.
Rows $0,2$ are then diagonal. Their first-column local
flag identity gives $\bar\chi_0=\bar\chi_2$, hence $y=1$
on the entire local decomposition group. Entry $12$
of the global representation has no two-step products
with this support. If it were a local coboundary in the
specified sense, odd-entry vanishing would kill it globally,
contrary to tightness of $B$. This proves the second
alternative of \texttt{LatShape}. All transformations
have been genuine integral matrices in a finite extension
of the original discretely valued fraction field, and
preserved the exact inertial labels. In particular no
specialization of an arbitrary pseudorepresentation has
been used.
\end{proof}

\begin{lemma}[Transport to every deeper tower]
\label{lem:rh-four-restr-q}
For the pair of Theorem~\ref{thm:rh-four-lat-shape} and
each original deeper fitted tower $T\ge T_0$, there is
a \texttt{Small} pair on $G_{T.F}$ satisfying the exact
\texttt{restrQ} conclusions and preserving all inertial
characters.
\end{lemma}
\begin{proof}
For each absolute place of $T.F$ over $\ell$ choose
$g$ transporting the fixed rational absolute place
to that place. The transported flag is the column flag
of $r(g)n$. In the residual support $nP$, row zero is
diagonal and the only possible entries of row one are
columns $0,1,2$. The flag of support $qQ$ has diagonal
one and zero entry $21$. Hence the first two leading
minors of the product are its first diagonal character
and the product of its first two diagonal characters.
They are units. The same Gaussian factorization as above,
now in the group of symplectic similitudes, gives the
lower symplectic unitriangular flag and upper triangular
similitude factor. For the latter factorization use
$g\mapsto J^{-1}g^{-t}J$ together with the multiplier;
uniqueness of the lower factor gives multiplier one
on it. Residually the lower factor has entry $21=0$
by elimination using row two, which has only columns
$0,2$. Thus it has support $qQ$. The conjugation identity
on the local decomposition group gives the original
ordered diagonal characters, as the characters on the
rational group are conjugation invariant.

Restriction preserves the global support and the killing
condition (here $\mathrm{Sb}=\varnothing$), and the new
mark is an $\ell$-power of the original diagonal mark.
If $s\in\mathrm{Mu}$, its image in the unramified quotient
$\widehat{\mathbb Z}$ has $\ell$-power torsion, hence is
zero. The cyclotomic principal-unit image of $s$ has
$\ell$-power torsion and is likewise zero. The residual
cyclotomic character, and the finite characters of
power-of-two order, are trivial on this local $\ell$-primary
abelian torsion. It follows that every original inertial
character is one on $s$. This verifies the exact
$\mathrm{Mu}$ condition. The transported flag entries
are integral in the same complete DVR, and restriction
preserves continuity. This proves all the claimed
\texttt{restrQ} properties, including \texttt{OverQ}
and the original local residual characters.
\end{proof}

The three new primes in the initial tower may be chosen
away from the discriminant of the fixed quadratic field
selected for the final trace bound. The further auxiliary
split primes of the residual-chain argument may also avoid
that finite set, by Chebotarev. This freedom will permit
the final fixed-field descent without changing that field.
